% Options for packages loaded elsewhere
% Options for packages loaded elsewhere
\PassOptionsToPackage{unicode}{hyperref}
\PassOptionsToPackage{hyphens}{url}
\PassOptionsToPackage{dvipsnames,svgnames,x11names}{xcolor}
%
\documentclass[
]{article}
\usepackage{xcolor}
\usepackage{amsmath,amssymb}
\setcounter{secnumdepth}{5}
\usepackage{iftex}
\ifPDFTeX
  \usepackage[T1]{fontenc}
  \usepackage[utf8]{inputenc}
  \usepackage{textcomp} % provide euro and other symbols
\else % if luatex or xetex
  \usepackage{unicode-math} % this also loads fontspec
  \defaultfontfeatures{Scale=MatchLowercase}
  \defaultfontfeatures[\rmfamily]{Ligatures=TeX,Scale=1}
\fi
\usepackage{lmodern}
\ifPDFTeX\else
  % xetex/luatex font selection
\fi
% Use upquote if available, for straight quotes in verbatim environments
\IfFileExists{upquote.sty}{\usepackage{upquote}}{}
\IfFileExists{microtype.sty}{% use microtype if available
  \usepackage[]{microtype}
  \UseMicrotypeSet[protrusion]{basicmath} % disable protrusion for tt fonts
}{}
\makeatletter
\@ifundefined{KOMAClassName}{% if non-KOMA class
  \IfFileExists{parskip.sty}{%
    \usepackage{parskip}
  }{% else
    \setlength{\parindent}{0pt}
    \setlength{\parskip}{6pt plus 2pt minus 1pt}}
}{% if KOMA class
  \KOMAoptions{parskip=half}}
\makeatother
% Make \paragraph and \subparagraph free-standing
\makeatletter
\ifx\paragraph\undefined\else
  \let\oldparagraph\paragraph
  \renewcommand{\paragraph}{
    \@ifstar
      \xxxParagraphStar
      \xxxParagraphNoStar
  }
  \newcommand{\xxxParagraphStar}[1]{\oldparagraph*{#1}\mbox{}}
  \newcommand{\xxxParagraphNoStar}[1]{\oldparagraph{#1}\mbox{}}
\fi
\ifx\subparagraph\undefined\else
  \let\oldsubparagraph\subparagraph
  \renewcommand{\subparagraph}{
    \@ifstar
      \xxxSubParagraphStar
      \xxxSubParagraphNoStar
  }
  \newcommand{\xxxSubParagraphStar}[1]{\oldsubparagraph*{#1}\mbox{}}
  \newcommand{\xxxSubParagraphNoStar}[1]{\oldsubparagraph{#1}\mbox{}}
\fi
\makeatother


\usepackage{longtable,booktabs,array}
\newcounter{none} % for unnumbered tables
\usepackage{calc} % for calculating minipage widths
% Correct order of tables after \paragraph or \subparagraph
\usepackage{etoolbox}
\makeatletter
\patchcmd\longtable{\par}{\if@noskipsec\mbox{}\fi\par}{}{}
\makeatother
% Allow footnotes in longtable head/foot
\IfFileExists{footnotehyper.sty}{\usepackage{footnotehyper}}{\usepackage{footnote}}
\makesavenoteenv{longtable}
\usepackage{graphicx}
\makeatletter
\newsavebox\pandoc@box
\newcommand*\pandocbounded[1]{% scales image to fit in text height/width
  \sbox\pandoc@box{#1}%
  \Gscale@div\@tempa{\textheight}{\dimexpr\ht\pandoc@box+\dp\pandoc@box\relax}%
  \Gscale@div\@tempb{\linewidth}{\wd\pandoc@box}%
  \ifdim\@tempb\p@<\@tempa\p@\let\@tempa\@tempb\fi% select the smaller of both
  \ifdim\@tempa\p@<\p@\scalebox{\@tempa}{\usebox\pandoc@box}%
  \else\usebox{\pandoc@box}%
  \fi%
}
% Set default figure placement to htbp
\def\fps@figure{htbp}
\makeatother


% definitions for citeproc citations
\NewDocumentCommand\citeproctext{}{}
\NewDocumentCommand\citeproc{mm}{%
  \begingroup\def\citeproctext{#2}\cite{#1}\endgroup}
\makeatletter
 % allow citations to break across lines
 \let\@cite@ofmt\@firstofone
 % avoid brackets around text for \cite:
 \def\@biblabel#1{}
 \def\@cite#1#2{{#1\if@tempswa , #2\fi}}
\makeatother
\newlength{\cslhangindent}
\setlength{\cslhangindent}{1.5em}
\newlength{\csllabelwidth}
\setlength{\csllabelwidth}{3em}
\newenvironment{CSLReferences}[2] % #1 hanging-indent, #2 entry-spacing
 {\begin{list}{}{%
  \setlength{\itemindent}{0pt}
  \setlength{\leftmargin}{0pt}
  \setlength{\parsep}{0pt}
  % turn on hanging indent if param 1 is 1
  \ifodd #1
   \setlength{\leftmargin}{\cslhangindent}
   \setlength{\itemindent}{-1\cslhangindent}
  \fi
  % set entry spacing
  \setlength{\itemsep}{#2\baselineskip}}}
 {\end{list}}
\usepackage{calc}
\newcommand{\CSLBlock}[1]{\hfill\break\parbox[t]{\linewidth}{\strut\ignorespaces#1\strut}}
\newcommand{\CSLLeftMargin}[1]{\parbox[t]{\csllabelwidth}{\strut#1\strut}}
\newcommand{\CSLRightInline}[1]{\parbox[t]{\linewidth - \csllabelwidth}{\strut#1\strut}}
\newcommand{\CSLIndent}[1]{\hspace{\cslhangindent}#1}



\setlength{\emergencystretch}{3em} % prevent overfull lines

\providecommand{\tightlist}{%
  \setlength{\itemsep}{0pt}\setlength{\parskip}{0pt}}



 


\makeatletter
\@ifpackageloaded{caption}{}{\usepackage{caption}}
\AtBeginDocument{%
\ifdefined\contentsname
  \renewcommand*\contentsname{Table of contents}
\else
  \newcommand\contentsname{Table of contents}
\fi
\ifdefined\listfigurename
  \renewcommand*\listfigurename{List of Figures}
\else
  \newcommand\listfigurename{List of Figures}
\fi
\ifdefined\listtablename
  \renewcommand*\listtablename{List of Tables}
\else
  \newcommand\listtablename{List of Tables}
\fi
\ifdefined\figurename
  \renewcommand*\figurename{Figure}
\else
  \newcommand\figurename{Figure}
\fi
\ifdefined\tablename
  \renewcommand*\tablename{Table}
\else
  \newcommand\tablename{Table}
\fi
}
\@ifpackageloaded{float}{}{\usepackage{float}}
\floatstyle{ruled}
\@ifundefined{c@chapter}{\newfloat{codelisting}{h}{lop}}{\newfloat{codelisting}{h}{lop}[chapter]}
\floatname{codelisting}{Listing}
\newcommand*\listoflistings{\listof{codelisting}{List of Listings}}
\makeatother
\makeatletter
\makeatother
\makeatletter
\@ifpackageloaded{caption}{}{\usepackage{caption}}
\@ifpackageloaded{subcaption}{}{\usepackage{subcaption}}
\makeatother
\usepackage{bookmark}
\IfFileExists{xurl.sty}{\usepackage{xurl}}{} % add URL line breaks if available
\urlstyle{same}
% fallback for those not using the hyperref driver hyperxmp:
\makeatletter
\@ifundefined{xmpquote}{\newcommand{\xmpquote}[1]{#1}}{}
\makeatother
\hypersetup{
  pdftitle={Dynamic Downsampling},
  pdfkeywords={\xmpquote{food webs}},
  colorlinks=true,
  linkcolor={blue},
  filecolor={Maroon},
  citecolor={Blue},
  urlcolor={Blue},
  pdfcreator={LaTeX via pandoc}}



\title{Dynamic Downsampling}
\author{Tanya Strydom %
%
\textsuperscript{%
%
1%
}%
; Andrew P. Beckerman %
%
\textsuperscript{%
%
1%
}%
}
\date{2026-07-31}

\usepackage{setspace}
\usepackage[left]{lineno}
\usepackage[letterpaper]{geometry}

\usepackage[nolists,noheads,markers]{endfloat}
\geometry{margin=2.5cm}

\begin{document}

\thispagestyle{empty}
{\bfseries\sffamily\Large Dynamic Downsampling}
\vfil
Tanya Strydom %
%
\textsuperscript{%
%
1%
}%
; Andrew P. Beckerman %
%
\textsuperscript{%
%
1%
}%

\vfil
{\small
\textbf{Abstract:} TODO
\vfil
\textbf{Keywords:} %
%
food webs%
}
\clearpage
\setcounter{page}{1}
\doublespacing
\linenumbers


\section{Introduction}\label{introduction}

Extinction simulation in networks are useful because they can inform us
as to the robustness {[}REF{]} or resilience {[}REF{]} of the system.
Understanding community level extinctions is of particular interest for
understanding community collapse and mass extinction events {[}REFs{]}.
However the nature of the fossil record means that we are only able to
do topological extinctions {[}REFs{]}

Dynamic extinctions (\emph{e.g.,} Curtsdotter et al., 2011) can be
useful because they are (probably) a better approximation of reality -
or at least not as `conservative'/best case scenario as what the
topological extinctions are (insert a brief description of how these
approaches are different). Something about how this is meaningful in
paleo but also real-world settings. Maybe also point out that typically
in dynamic context we are using species agnostic models \emph{e.g.,}
niche and atn which lack the real world tractability that we want.

However in order to \emph{correctly} match the philosophy of these
dynamic networks we need a way in which to turn empirical metawebs into
networks that are more representative of realised networks, \emph{i.e.,}
there is some form of link distribution constraint (Strydom, Dunhill, et
al., 2026). Although there has been some work on downsampling of paleo
webs (Roopnarine, 2006) we lack any ground truthing of what the
implications are of embedding these ecological assumptions. Additionally
this downsampling approach makes insidious assumptions as to what is
driving the distribution of links in a network. Possibly also add a
sentence here about how the other work has focused on the niche model
(Curtsdotter et al., 2011; Jonsson et al., 2015) which itself has its
own sets of ecological assumptions/embeddings.

Here we present a selection of downsampling approach that are built on
different assumptions as to link constraints. Using these downsampling
approaches we recreate the Curtsdotter et al. (2011) approach and test
if we see difference in topological and dynamic extinctions, as well as
if the different input networks have different signals. We do this using
the same dataset from Karapunar et al. (2026).

\section{Methods}\label{methods}

\#\#~Data

Seven locations (Karapunar et al., 2026). Here we use the guild-level
community data to keep the size of the networks manageable and well
within the bounds used by Curtsdotter et al. (2011). Additionally we
added plankton and zooplankton data from EcoGenie models {[}REF{]} as a
way to enlarge the `base' of the network (again to better represent a
realised network).

\subsection{Network construction}\label{network-construction}

For the empirical networks we used the PFWIM model (Shaw et al., 2024).
\emph{Insert brief description about trait assignments and mechanistic
assumptions, mention we wanted to go as broad as possible so we used
fully categorical rules}. This was then used as the base metaweb for
each community. Each metaweb was then downsampled to a target
connectance (see below for the different downsampling approaches).
Additionally the target connectance and species richness was used to
parameterise a niche model network (Williams \& Martinez, 2000). We also
constructed ATN networks (Brose et al., 2006).

\subsection{Downsampling approaches}\label{downsampling-approaches}

Maybe just a brief interlude about how a metaweb represents all feasible
interactions (so capturing the entire diet breadth of a specific
species). Idea here is to prime that this is all links but acknowledge
that in reality species are rarely feeding across their entire diet
breadth.

\emph{Niche:} Attempt to create a network that appears similar to a
standard niche model network this is done by estimating a
one-dimensional niche axis from the interaction matrix using the leading
singular vectors of the Singular Value Decomposition (SVD). Species are
then assigned positions along a common latent niche dimension, which is
subsequently used to estimate consumer niche centres and niche breadths.

\emph{Power-law (Roopnarine, 2006):} Based around ideas of extinctions
and which links/specie will be lost - with the idea that generalist
species will be the most robust and as such will retain their links.
Consumers with larger trophic breadth (higher generality) receive larger
retention probabilities and are therefore less likely to lose
interactions.

\emph{Degree-product:} Product of the out-degreee (generality) of the
predator and the in-degree (vulnerability) of the consumer are used to
assign a probability for each link. Links are then randomly removed
based on that probability (note to self check if it is iterative each
time so that the degrees are re-evaluated, it think this is the case).
Here the ecological assumption is that generalist predators and
vulnerable prey are probably not feeding/being fed on by everything,
whereas for specialist species those links are probably occurring thus
it makes sense to prune links from the most connected species pairs

\emph{Random:} Links randomly removed until target connectance is
reached, with the exception if a species becomes isolated in which case
a different link is target for removal. This is the `null model'

\subsection{Simulations}\label{simulations}

For each community we create a network, sample a body size based on the
categorical size classes, this is a truncated normal distribution. Body
sizes are used to build the ATN (mention here that we vary the threshold
value between 0.01-0.12 and select the network that has the closest
connectance to the target one) We also select a target connectance
between 0.05-0.15. Each network is then downsampled to the target
connectance. This represents the `creation' stage. At this point we run
a set of topological cascading extinctions for the difference extinction
scenarios. Using the END (Delmas et al., 2017) each network is then run
through a `burn-in' period for 5000 time steps or until it has reached a
steady state. Here we run another set of topological extinctions. For
the dynamic extinctions primary extinctions are determined based on the
rules of the different extinction scenarios and secondary extinctions
occur when the biomass of a species falls below the threshold value of
1e-12. After every primary extinction event (where the biomass of the
target species is set to 1e-12) we run 5000 timesteps before the next
extinction event to allow for biomass simulations to occur.

\subsection{Analyses}\label{analyses}

For each extinction run we calculate the R50 robustness (Jonsson et al.,
2015) as well as the primary and secondary extinction rates. Add breif
description of how robustness is calculated. TBD the breakdown of the
statistical approaches used to understand the effect of downsampling
approach/network type on inferred network robustness and extinction
dynamics.

\subsection{Software}\label{software}

Network construction and extinction simulations were run in Julia
v1.12.6 (Bezanson et al., 2017), we used
\texttt{EcologicalNetworksDynamics.jl} (Lajaaiti et al., 2025) for
dynamic simulations, \texttt{PFIM.jl} {[}REF{]} for construction of
metawebs, \texttt{FoodWebTools.jl} {[}REF{]} for network downsampling,
and Extinctions.jl {[}REF{]} for extinctions and robustness
calculations. Statistical analyses were run in R v4.5.2 {[}REF{]}. All
code can be found at \emph{TODO}

\begin{quote}
Keep the narrative extinction centric for now but remember we also have
the topology of each network before the various extinction simulations
begin
\end{quote}

\section{Downsampling approach results in different network
structure}\label{downsampling-approach-results-in-different-network-structure}

Here we can actually look at the topology and see if the different
networks actually look different in multivariate space?

\section{Differences in inferred robustness of topological and dynamic
unaffected by network
type}\label{differences-in-inferred-robustness-of-topological-and-dynamic-unaffected-by-network-type}

Still see that topological robustness is more conservative than dynamic
extinctions

\section{Network types result in differences in dynamic
extinctions}\label{network-types-result-in-differences-in-dynamic-extinctions}

Report here about the secondary extinction rates as well as the pairwise
contrasts.

\#\#~Relative to the niche model, niche downsampling is the closest
match

Most often the ones to not be different from each other. Worth pointing
out that the ATN behaves very differently as well. So again we are back
at are we modelling network or model behaviour (Strydom, Karapunar, et
al., 2026)

\section{Conclusion}\label{conclusion}

\section*{References}\label{references}
\addcontentsline{toc}{section}{References}

\protect\phantomsection\label{refs}
\begin{CSLReferences}{1}{0}
\bibitem[\citeproctext]{ref-bezansonJuliaFreshApproach2017}
Bezanson, J., Edelman, A., Karpinski, S., \& Shah, V. (2017). Julia: {A
Fresh Approach} to {Numerical Computing}. \emph{SIAM Review},
\emph{59}(1), 65--98. \url{https://doi.org/10.1137/141000671}

\bibitem[\citeproctext]{ref-broseConsumerResourceBodySize2006}
Brose, U., Jonsson, T., Berlow, E. L., Warren, P., Banasek-Richter, C.,
Bersier, L.-F., Blanchard, J. L., Brey, T., Carpenter, S. R.,
Blandenier, M.-F. C., Cushing, L., Dawah, H. A., Dell, T., Edwards, F.,
Harper-Smith, S., Jacob, U., Ledger, M. E., Martinez, N. D., Memmott,
J., \ldots{} Cohen, J. E. (2006). Consumer--{Resource Body-Size
Relationships} in {Natural Food Webs}. \emph{Ecology}, \emph{87}(10),
2411--2417.
\url{https://doi.org/10.1890/0012-9658(2006)87\%5B2411:CBRINF\%5D2.0.CO;2}

\bibitem[\citeproctext]{ref-curtsdotterRobustnessSecondaryExtinctions2011}
Curtsdotter, A., Binzer, A., Brose, U., De Castro, F., Ebenman, B.,
Eklöf, A., Riede, J. O., Thierry, A., \& Rall, B. C. (2011). Robustness
to secondary extinctions: {Comparing} trait-based sequential deletions
in static and dynamic food webs. \emph{Basic and Applied Ecology},
\emph{12}(7), 571--580. \url{https://doi.org/10.1016/j.baae.2011.09.008}

\bibitem[\citeproctext]{ref-delmasSimulationsBiomassDynamics2017}
Delmas, E., Brose, U., Gravel, D., Stouffer, D. B., \& Poisot, T.
(2017). Simulations of biomass dynamics in community food webs.
\emph{Methods in Ecology and Evolution}, \emph{8}(7), 881--886.
\url{https://doi.org/10.1111/2041-210X.12713}

\bibitem[\citeproctext]{ref-jonssonReliabilityR50Measure2015}
Jonsson, T., Berg, S., Pimenov, A., Palmer, C., \& Emmerson, M. (2015).
The reliability of {R50} as a measure of vulnerability of food webs to
sequential species deletions. \emph{Oikos}, \emph{124}(4), 446--457.
\url{https://doi.org/10.1111/oik.01588}

\bibitem[\citeproctext]{ref-karapunarNoGlobalCollapse2026}
Karapunar, B., Strydom, T., Beckerman, A. P., Ridgwell, A., Wignall, P.
B., Dunne, J. A., Little, C. T. S., Hull, P., Pimiento, C., \& Dunhill,
A. (2026, February 25). \emph{No global collapse of food webs across the
{Permian-Triassic Mass Extinction}}.
\url{https://doi.org/10.64898/2026.02.24.707709} (Pre-published)

\bibitem[\citeproctext]{ref-lajaaitiEcologicalNetworksDynamicsjlJuliaPackage2025}
Lajaaiti, I., Bonnici, I., Kéfi, S., Mayall, H., Danet, A., Beckerman,
A. P., Malpas, T., \& Delmas, E. (2025).
{EcologicalNetworksDynamics}.jl: {A Julia} package to simulate the
temporal dynamics of complex ecological networks. \emph{Methods in
Ecology and Evolution}, \emph{16}(3), 520--529.
\url{https://doi.org/10.1111/2041-210X.14497}

\bibitem[\citeproctext]{ref-roopnarineExtinctionCascadesCatastrophe2006}
Roopnarine, P. D. (2006). Extinction {Cascades} and {Catastrophe} in
{Ancient Food Webs}. \emph{Paleobiology}, \emph{32}(1), 1--19.
\url{https://www.jstor.org/stable/4096814}

\bibitem[\citeproctext]{ref-shawFrameworkReconstructingAncient2024}
Shaw, J. O., Dunhill, A. M., Beckerman, A. P., Dunne, J. A., \& Hull, P.
M. (2024, January 30). \emph{A framework for reconstructing ancient food
webs using functional trait data}.
\url{https://doi.org/10.1101/2024.01.30.578036} (Pre-published)

\bibitem[\citeproctext]{ref-strydomScalingMetawebsRealised2026}
Strydom, T., Dunhill, A. M., Dunne, J. A., Poisot, T., \& Beckerman, A.
P. (2026). \emph{From {Metawebs} to {Realised Webs}: {A Framework} for
{Ecological Network Representation} under {Global Change}}.
\url{https://ecoevorxiv.org/repository/view/11514/}

\bibitem[\citeproctext]{ref-strydomAreWeMapping2026}
Strydom, T., Karapunar, B., Dunne, J. A., Hull, P., Little, C. T. S.,
Pimiento, C., Ridgwell, A., Beckerman, A. P., \& Dunhill, A. M. (2026).
\emph{Are {We Mapping Ecosystems} or {Models}? {Framework Choices
Dominate Food Web Topology} and {Extinction Inferences}}.
\url{https://ecoevorxiv.org/repository/view/13495/}

\bibitem[\citeproctext]{ref-Williams2000Simple}
Williams, R. J., \& Martinez, N. D. (2000). Simple rules yield complex
food webs. \emph{Nature}, \emph{404}(6774), 180--183.
\url{https://doi.org/10.1038/35004572}

\end{CSLReferences}





\end{document}
